\documentclass[11pt,a4paper]{article}
\usepackage[T1]{fontenc}
\usepackage{isabelle,isabellesym}
\usepackage{amssymb,amsmath}
\usepackage{pdfsetup}
\hypersetup{pdftitle={Dependent Families and Displayed Algebras over Structural HIT Syntax}, pdfauthor={Arthur Freitas Ramos, David Barros Hulak, Ruy de Queiroz}}
\urlstyle{rm}
\isabellestyle{it}
\title{Dependent Families and Displayed Algebras over Structural HIT Syntax}
\author{Arthur Freitas Ramos \and David Barros Hulak \and Ruy de Queiroz}
\begin{document}
\maketitle
\begin{abstract}
For set-valued families over structural higher-inductive-type signatures,
we prove a local-to-global coherence theorem: under explicit sufficient
interchange laws, positive elementary recursive-slot naturality is equivalent
to compatibility with all canonical tuple naturality cells, on every fiber
element and under every valid constructor shell. Displayed point methods
construct dependent folds whose coherence descends to indexed path quotients
under uniform boundary and contextual laws.
These constructions use the published structural HIT syntax and support
external parameters and substitution of whole displayed algebras. Examples
separate coherence of a selected section from all-fiber compatibility.
The result is an external dependent semantic layer, not a dependent object
calculus or a general model of higher inductive types.
\end{abstract}

\section{Main theorem: from elementary squares to coherent transport}

A path constructor is a family of generating paths
\(\ell_{\vec x}:L(\vec x)\longrightarrow R(\vec x)\), where its two
endpoints are expressions in the imported signature. A dependent family
assigns a set \(F(t)\) to each base term and mutually inverse transport maps
to each admitted generating edge. The question is whether this transport
respects the equations imposed on paths, for \emph{every} value in the source
fiber, rather than only for the values of one chosen section.

\paragraph{The common hypotheses.}
Fix a well-formed original signature, external parameter carriers and the
admitted constructor-tree grammar. Same-sort parameters are excluded in this
fragment; external parameters remain fixed during argument motion. Constructor
contexts may contain a recursive hole, and a valid outer \emph{shell} is such
a context surrounding a path. Assume the sufficient interchange interface:
action in either of two distinct recursive slots is unchanged, as a total HOL
function, when the other slot changes its endpoint; the two initial slot
actions commute on every value of the appropriate fiber. These laws hold
uniformly under every valid outer shell. They are strong sufficient
hypotheses, not a characterization of all possible dependent families.

\paragraph{Two naturality conditions.}
An \emph{elementary move} changes one recursive argument along one positive
admitted contextual edge, leaving the other arguments fixed.
\emph{Elementary naturality} requires its square with each path constructor
to commute under every valid shell and on every source-fiber value. The
admitted edge may already lie inside a constructor context: elementary does
not mean only an unwrapped original generator.

A \emph{tuple motion} instead assigns an arbitrary finite signed path to each
recursive argument. Original endpoint expressions act on these paths to give
\(L(\alpha)\) and \(R(\alpha)\). \emph{Tuple naturality} requires the two
parallel composite paths
\[
 L(\alpha)\,;\,\ell_{\vec y}
 \qquad\hbox{and}\qquad
 \ell_{\vec x}\,;\,R(\alpha)
\]
to induce the same transport, again on every source-fiber value and under
every valid shell. Here a semicolon means traversal of the left path followed
by the right path.

\paragraph{Informal main theorem.}
\emph{Under the common hypotheses above, elementary naturality is equivalent
to tuple naturality.} The exact exported Isabelle theorem is
\texttt{local\_global\_signature\_coherence}; it is stated outside a locale
as the equivalence of \texttt{locally\_natural\_signature\_family} and
\texttt{coherent\_signature\_family}. Both predicates include the same
signature, admitted-family and sufficient interchange requirements. Thus the
equivalence does not erase those requirements. Its substantive direction
extends elementary squares along signed paths, serializes tuple motions, and
uses proved endpoint functoriality and interchange to compare the two endpoint
actions. The reverse direction specializes a tuple motion to one elementary
move.

\paragraph{Why this is useful.}
To obtain an action of the combined interchange/naturality quotient, one can
verify elementary naturality rather than separately prove compatibility for
every tuple of composite and reversed argument paths. This reduces the form
of the obligations, not their quantification: all valid shells and all fiber
values still matter, and no finite or decidable checking procedure is claimed.
The derived quotient transport is independent of the representative of an
endpoint-indexed path class.

A separate contextual displayed algebra supplies typed point methods,
base-generator boundary coherence for the recursively defined fold, and
uniform contextual method naturality for arbitrary correctly typed recursive
results. Earlier displayed-HIT boundary methods provide one sufficient way
to establish the base-generator condition. The point methods recursively
define a fold \(f\); the construction does not assume a global section. If \([p]:x\longrightarrow y\)
is in the corresponding quotient hom set and \(x\) is in the admitted tree
grammar, \texttt{constructed\_coherent\_elimination} proves
\[
 \operatorname{transport}_{[p]}(f(x))=f(y)
 \qquad\hbox{and}\qquad f(y)\in F(y).
\]
This is the constructed-fold consequence of the main coherence result, not
a general dependent object-language eliminator or a typed HIT universal
property. The following sections explain its interfaces and limitations.

\section{Scope and relationship to the imported calculus}
The published metatheory~\cite{HITMetatheory} has nondependent function spaces
and constant-motive eliminators. We import its actual signatures, syntax,
endpoint expressions and substitution theory; we do not relabel its type
system as dependent. The new layer consists of fibers that are subsets of a
common HOL carrier, indexed by base vertices, with bijective transport along
admitted edges. The admitted-edge restriction permits empty graphs and keeps
invalid constructor instances out of the path grammar.

Dependent sums and products allow the second family to depend on both a base
vertex and an element of the first family. Functions representing dependent
products are normalized outside the source fiber. A Boolean fiber with
negation around a loop illustrates nontrivial monodromy and the absence of
sections. The connected-section classification assumes explicit connecting
paths from a root and must only be used on a domain satisfying that hypothesis.

\section{Constructed sections and substitution}
Generator coherence of a supplied assignment extends to paths, but is not
alone a construction of an eliminator. The displayed-algebra modules instead
recursively construct folds from point methods on the original constructor
syntax, prove their typing and uniqueness on the designated grammar, and
interpret original endpoint expressions. Mixed constructors provide
induction hypotheses only at \texttt{Rec} slots. Opaque \texttt{Par} arguments
belong to specified external carriers; same-sort parameters of the HIT being
eliminated are excluded. A path-constructor arity containing such a parameter
has no admitted instances in this fragment. Graph-path typing asserts
admissibility and endpoint matching; typing in the imported object calculus
additionally requires the explicit signature and external-parameter typing
hypotheses. Arbitrary parameter carriers are not assumed well typed.
Boundary equations quantify arbitrary correctly
typed recursive results, not only the values of the completed fold.

Recursive-slot contexts extend this construction through point constructors
under uniform local displayed-method naturality. External parameters remain
fixed along these contextual paths. Substitution commutes with the constructed
fold and pulls back whole displayed and contextual algebras, with substituted
parameter carriers, ordinary arguments and contexts. The source admissibility
premises are retained: substitution need not reflect the designated grammar.

\section{Equations, quotients and all-fiber coherence}
The typed path presentation closes generating equations under equivalence,
composition, reversal and inverse cancellation. Its quotients are actual
sets of equivalent paths in each indexed hom-set, with identity, composition
and inverse operations satisfying groupoid laws. These are indexed hom-sets,
not a globally tagged arrow datatype: class sets at different endpoint pairs
can overlap. A chosen representative is used only on genuine nonempty
classes, and the operations and transport are proved independent of that
choice. This is an abstract HOL construction, not an executable normalizer
or a decision procedure for quotient equality.

A family descends to the quotient when every generating equation acts equally
on \emph{all} elements of its source fiber. This is stronger than coherence
of one constructed section. The fold descends only from starting terms in
its constructor grammar; the existence of an empty path at an arbitrary raw
term does not extend that grammar. Presented path maps send edges to paths
and, when generating equations are preserved, induce maps of hom-set
quotients preserving identities, composition and inverses.

Independent recursive slots satisfy interchange under explicit sufficient
transport laws. The current slot-independence hypotheses are equalities of
total transport functions, whereas their commutation law is restricted to
the appropriate fiber. Positive squares yield all signed cases and grids
for arbitrary paths in two distinct recursive slots. A non-singleton
integer-vector example satisfies these laws. A counterexample fixes a
coherent section while violating interchange on another fiber element;
section coherence alone cannot justify the all-fiber quotient.
A separate count-affine model has infinite fibers and realizes an actual
recursive-generator naturality square with arbitrary residual integer values.
Its transport depends only on endpoints, so it witnesses consistency rather
than faithfulness of the specified path equations.

\section{Original endpoints and generator naturality}
The endpoint action recursively follows the imported \texttt{EArg} and
\texttt{ECon} syntax. Canonical path-constructor instances retain their
constructor identifier and argument tuple. Endpoint motion has proved typed
boundaries, and an explicit argument graph serializes simultaneous recursive
argument motions into elementary moves while holding external parameters
fixed. Beyond typing, the endpoint action now has verified identity,
composition and inversion laws modulo interchange. Its value on a tuple is
proved equivalent to the path image of the tuple's elementary serialization.

A combined presentation adds canonical generator-naturality cells under
valid outer constructor shells to interchange. Within the same sufficient
interchange interface, compatibility with all tuple naturality cells is
proved equivalent to positive elementary recursive-slot naturality. Both
conditions range over every valid outer shell and every source-fiber value.
Elementary edges may already be contextual, not only bare generators.
Accordingly these are conditional coherence results, not a claim that
arbitrary point methods or endpoint actions automatically define a coherent
HIT algebra.

For a contextual displayed algebra with the same methods and transport,
the recursively constructed fold descends to this combined coherence
quotient. The theorem includes typing of its value at the target and retains
the constructor-grammar premise at the source. It assumes no pre-existing
global section.

\section{A combined path-sensitive test model}

The theory \texttt{DHIT\_Path\_Sensitive\_Example} uses one original signature
with a base point \(z\), a unary constructor \(u\), a binary constructor \(b\),
a recursive path constructor \(r(x):u(x)\longrightarrow x\), and a nullary
loop \(\ell:z\longrightarrow z\). Every fiber is the set of integers.
Transport along the loop negates the integer, while transport along a collapse
generator is the identity. Constructor contexts preserve this underlying
generator action.

The interpretation satisfies the full local naturality interface and hence
the combined coherence interface. In particular, the recursive naturality
square compares \(u(\ell);r(z)\) with \(r(z);\ell\); both act by negation.
In the binary constructor \(b(z,z)\), the loop in the left slot and the loop
in the right slot may be traversed in either order. These are distinct path
words with nonidentity individual actions, and both composites act identically
on every integer. Thus the same model exercises recursive-generator
naturality and two-slot interchange.

This example is not endpoint-determined: \(\ell\) and the empty path have
the same endpoints but act differently on \(1\). All-fiber quotient descent
therefore proves that they cannot be equal in the combined path presentation.
The fibers here are constant, the two slot actions are the same commuting
involution, and their two-step action is the identity. No faithfulness of the
entire presentation or displayed-fold interpretation of this particular model
is claimed; the earlier examples provide the constructed-fold instances.

\section{Relationship to prior work}
Transported families, displayed algebras and dependent elimination are
established mathematics; see the HoTT book~\cite{HoTTBook2013}, the groupoid
interpretation~\cite{HofmannStreicher1998}, homotopy-initial
algebras~\cite{Sojakova2015}, signatures and induction
principles~\cite{KaposiKovacs2020}, and groupoid constructions of
HITs~\cite{DybjerMoeneclaey2018,VeltriVanDerWeide2021}. The contribution here is
the scoped Isabelle/HOL development and its proved connection to the imported
structural syntax, rather than a new general theory of dependent elimination.
The AFP entries Category3~\cite{Category3AFP} and Catoids~\cite{CatoidsAFP}
provide related established library infrastructure. They are cited for
context; this entry does not import or reuse their formal theories.

\section{Boundaries}
No dependent lambda-calculus syntax, universe, univalence principle, Kan
operation, groupoid-valued motive, general higher-dimensional HIT, or full
typed universal property is established. The equations added here are not
claimed complete for the equality of the original calculus. Path cancellation
subject reduction and termination do not establish normalization of a
dependent calculus; no normalization theorem is transferred from the imported
nondependent development.

\tableofcontents
\input{session}
\bibliographystyle{abbrv}
\bibliography{root}
\end{document}
